%--------------------------------------------------------------------------------------
% Page layout setup
%--------------------------------------------------------------------------------------
% we need to redefine the pagestyle plain
% another possibility is to use the body of this command without \fancypagestyle
% and use \pagestyle{fancy} but in that case the special pages
% (like the ToC, the References, and the Chapter pages)remain in plane style

\pagestyle{plain}

\setcounter{tocdepth}{3}
%\sectionfont{\large\upshape\bfseries}
\setcounter{secnumdepth}{3}

\sloppy % Margón túllógó sorok tiltása.
\widowpenalty=10000 \clubpenalty=10000 %A fattyú- és árvasorok elkerülése
\def\hyph{-\penalty0\hskip0pt\relax} % Kötőjeles szavak elválasztásának engedélyezése


%--------------------------------------------------------------------------------------
% Setup hyperref package
%--------------------------------------------------------------------------------------
\hypersetup{
    % bookmarks=true,            % show bookmarks bar?
    unicode=true,              % non-Latin characters in Acrobat's bookmarks
    pdftitle={\vikcim},        % title
    pdfauthor={\szerzoMeta},    % author
    pdfsubject={\vikdoktipus}, % subject of the document
    pdfcreator={\szerzoMeta},   % creator of the document
    pdfproducer={},    % producer of the document
    pdfkeywords={},    % list of keywords (separate then by comma)
    pdfnewwindow=true,         % links in new window
    colorlinks=true,           % false: boxed links; true: colored links
    linkcolor=black,           % color of internal links
    citecolor=black,           % color of links to bibliography
    filecolor=black,           % color of file links
    urlcolor=black             % color of external links
}


\input{include/code-languages} % syntax-highlighting colors, minted setup, Refinery/Tamarin/Lean

%--------------------------------------------------------------------------------------
% Set up listings
%--------------------------------------------------------------------------------------
\definecolor{lightgray}{rgb}{0.95,0.95,0.95}
\lstset{
	basicstyle=\scriptsize\ttfamily, % print whole listing small
	keywordstyle=\color{codeKeyword}\bfseries,
	identifierstyle=, % nothing happens
	commentstyle=\color{codeComment}\itshape,
	stringstyle=\color{codeString}\scriptsize,
	showstringspaces=false, % no special string spaces
	aboveskip=3pt,
	belowskip=3pt,
	backgroundcolor=\color{lightgray},
	columns=flexible,
	keepspaces=true,
	escapeinside={(*@}{@*)},
	captionpos=b,
	breaklines=true,
	frame=single,
	float=!ht,
	tabsize=2,
	literate=*
		{á}{{\'a}}1	{é}{{\'e}}1	{í}{{\'i}}1	{ó}{{\'o}}1	{ö}{{\"o}}1	{ő}{{\H{o}}}1	{ú}{{\'u}}1	{ü}{{\"u}}1	{ű}{{\H{u}}}1
		{Á}{{\'A}}1	{É}{{\'E}}1	{Í}{{\'I}}1	{Ó}{{\'O}}1	{Ö}{{\"O}}1	{Ő}{{\H{O}}}1	{Ú}{{\'U}}1	{Ü}{{\"U}}1	{Ű}{{\H{U}}}1
}


%--------------------------------------------------------------------------------------
% Set up theorem-like environments
%--------------------------------------------------------------------------------------
% Using ntheorem package -- see http://www.math.washington.edu/tex-archive/macros/latex/contrib/ntheorem/ntheorem.pdf

\theoremstyle{plain}
\theoremseparator{.}
\newtheorem{example}{\pelda}

\theoremseparator{.}
%\theoremprework{\bigskip\hrule\medskip}
%\theorempostwork{\hrule\bigskip}
\theorembodyfont{\upshape}
\theoremsymbol{{\large \ensuremath{\cdot}}}
\newtheorem{definition}{\definicio}

\theoremseparator{.}
%\theoremprework{\bigskip\hrule\medskip}
%\theorempostwork{\hrule\bigskip}
\newtheorem{theorem}{\tetel}

% More theorem-like environments, numbered together with theorems (Lemma 3, Theorem 4, ...).
% Italic body: formal statements.
\theorembodyfont{\itshape}
\theoremsymbol{}
\newtheorem{lemma}[theorem]{Lemma}
\newtheorem{proposition}[theorem]{Proposition}
\newtheorem{corollary}[theorem]{Corollary}
\newtheorem{invariant}[theorem]{Invariant}
\newtheorem{property}[theorem]{Property}
\newtheorem{claim}[theorem]{Claim}
% Upright body: commentary.
\theorembodyfont{\upshape}
\newtheorem{remark}[theorem]{Remark}
\newtheorem{notation}[theorem]{Notation}
% Unnumbered proof that ends with a box.
\theoremstyle{nonumberplain}
\theoremheaderfont{\itshape}
\theoremsymbol{\ensuremath{\square}}
\newtheorem{proof}{Proof}
\theoremstyle{plain}
\theoremheaderfont{\normalfont\bfseries}
\theoremsymbol{}


%--------------------------------------------------------------------------------------
% Some new commands and declarations
%--------------------------------------------------------------------------------------
\newcommand{\code}[1]{{\upshape\ttfamily\scriptsize\indent #1}}
\newcommand{\doi}[1]{DOI: \href{https://doi.org/\detokenize{#1}}{\raggedright{\texttt{\detokenize{#1}}}}} % A hivatkozások közt így könnyebb DOI-t megadni.

% English abbreviations; include/thesis-en.tex defines the full set, these keep
% Hungarian builds compiling when a text uses them.
\providecommand{\eg}{e.g.\@\xspace}
\providecommand{\ie}{i.e.\@\xspace}
\providecommand{\etal}{et al.\@\xspace}

\DeclareMathOperator*{\argmax}{arg\,max}
%\DeclareMathOperator*[1]{\floor}{arg\,max}
\DeclareMathOperator{\sign}{sgn}
\DeclareMathOperator{\rot}{rot}


%--------------------------------------------------------------------------------------
% Setup captions
%--------------------------------------------------------------------------------------
\captionsetup[figure]{
	width=.75\textwidth,
	aboveskip=10pt}

\renewcommand{\captionlabelfont}{\bf}

% Number column in the list of figures and tables, wide enough for A.3.1.
\makeatletter
\renewcommand*{\l@figure}{\@dottedtocline{1}{1.5em}{3.2em}}
\renewcommand*{\l@table}{\@dottedtocline{1}{1.5em}{3.2em}}
% The same for subsections in the table of contents, for numbers like A.10.10.
\renewcommand*{\l@subsection}{\@dottedtocline{2}{1.5em}{3.5em}}
\renewcommand*{\l@subsubsection}{\@dottedtocline{3}{5em}{4.1em}}
\makeatother

% Ragged right bibliography: long URLs and DOIs otherwise stretch the line.
\renewcommand{\bibfont}{\raggedright}
% The bibliography's line in the contents and its bookmark come from its own heading, so both point
% at the page it starts on (an \addcontentsline before \bibliography points at the page before).
\renewcommand{\bibsection}{\chapter*{\bibname}\addcontentsline{toc}{chapter}{\bibname}}
%\renewcommand{\captionfont}{\footnotesize\it}

%--------------------------------------------------------------------------------------
% Hyphenation exceptions
%--------------------------------------------------------------------------------------
\hyphenation{Shakes-peare Mar-seilles ár-víz-tű-rő tü-kör-fú-ró-gép}


\author{\vikszerzo}
\title{\viktitle}

%--------------------------------------------------------------------------------------
% Todo notes
%--------------------------------------------------------------------------------------
% The page margins are 15mm (see \marginsize in draft.tex): too narrow for a readable note
% that does not cover the text or run off the bottom of the page. So every \todo prints
% inline, a block across the text width under the line it is written on. \todo[noinline]{...}
% still puts a word or two in the margin, in a box about 10mm wide, clear of the text and of the
% page edge; anything longer than "cite" or "typo" belongs inline.
% todonotes checks for at least 2cm at \begin{document}; the narrow width is deliberate,
% so it is set after that check. Both kinds disappear with todonotes' disable option.
\setlength{\marginparwidth}{2cm}
\AtBeginDocument{\setlength{\marginparwidth}{11mm}\setlength{\marginparsep}{1.5mm}}
\makeatletter
\@ifpackageloaded{todonotes}{\setlength{\@todonotes@textwidth}{9.5mm}\presetkeys{todonotes}{inline}{}}{}
\makeatother
\newcommand{\todoinline}[1]{\todo[inline]{#1}}

% Two kinds of notes: orange \todo for the student; blue \advisortodo for the advisors
% (a draft to confirm, or a link to fill in). Both disappear with todonotes' disable option.
% A beat: the plan of one paragraph (see the Writing chapter of the guide). Prints as a student todo.
%   \beat{what it establishes}{its evidence}{what the running example shows here}{what it leads to}
\newcommand{\beat}[4]{\todo[inline, caption={Beat: #1}]{\textbf{Beat.} \textit{Establishes:} #1. \textit{Evidence:} #2. \textit{Running example:} #3. \textit{Leads to:} #4.}}
% A heading's pointer to its instruction and example in the guide chapter "The paper, chapter by
% chapter" (guide/05-chapter-by-chapter.tex). Prints as a student todo, so the final build drops it.
%   \guidetodo{Context}{app:chapters:context}
\newcommand{\guidetodo}[2]{\todo[inline, caption={#1}]{\textbf{#1.} What goes here, with an example: \cref{#2}.}}
% The guide prints its beats as text, so they stay out of the paper's list of todos, and marks a
% heading whose section holds something a skimming reviewer reads.
% (\leavevmode before \color: a colour switch between a heading and its text would let the page
% break there and leave the heading alone at the bottom of a page.)
\newcommand{\guidebeat}[4]{\par{\small\leavevmode\color{gray!30!black}\textbf{Beat.} \textit{Establishes:} #1. \textit{Evidence:} #2. \textit{Running example:} #3. \textit{Leads to:} #4.\par}\smallskip}
\newcommand{\skimmark}{\enspace{\small\faEye}}

% The running example, named once in draft.tex (\exname) and revisited in every chapter.
%   \begin{runningexample}[what this chapter adds] ... \end{runningexample}
\newtcolorbox{runningexample}[1][]{enhanced, breakable, sharp corners, frame hidden,
  colback=teal!5, borderline west={2pt}{0pt}{teal!60!black}, left=6pt, right=4pt, top=3pt, bottom=3pt,
  before upper={{\raggedright\textbf{Running example, \exname\if\relax\detokenize{#1}\relax\else: #1\fi}\par}\smallskip}}

% Long quotations: mandatory attribution argument, including a citation and locator.
%   \begin{sourcequote}{Author, \cite[p.~12]{key}} ... \end{sourcequote}
\newtcolorbox{sourcequote}[1]{enhanced, breakable, sharp corners, frame hidden,
  colback=white, borderline west={2pt}{0pt}{gray!60},
  left=8pt, right=4pt, top=4pt, bottom=4pt,
  after upper={\par\nobreak\smallskip{\raggedleft\normalfont #1\par}}}
% The final build (final.tex) leaves the guide out, so the paper names its sections instead
% of referencing them, and the guide names the paper's files.
\newcommand{\guideref}[1]{\enquote{#1} in the guide}
\makeatletter
% File paths break after / and - (\path from the url package), so they fit narrow table columns.
\g@addto@macro\UrlBreaks{\do\-}
\newcommand{\deftplref}[4]{\@namedef{tpl@#1}{#2 (\path{content/#4.tex})}\@namedef{Tpl@#1}{#3 (\path{content/#4.tex})}}
% An unknown key prints ?? and warns, as an undefined \ref does, instead of printing nothing.
\newcommand{\tdk@tplref}[2]{\@ifundefined{#1@#2}{\PackageWarning{tdk}{\string\tplref\space key `#2' undefined; add it with \string\deftplref}\textbf{??}}{\@nameuse{#1@#2}}}
\newcommand{\tplref}[1]{\tdk@tplref{tpl}{#1}}
\newcommand{\Tplref}[1]{\tdk@tplref{Tpl}{#1}}
\makeatother
\deftplref{ch:intro}{the introduction}{The introduction}{01-introduction}
\deftplref{ch:conclusion}{the conclusion}{The conclusion}{07-conclusion}
\deftplref{ch:approach}{the approach chapter}{The approach chapter}{03-approach}
\deftplref{ch:implementation}{the implementation chapter}{The implementation chapter}{04-implementation}
\deftplref{ch:evaluation}{the evaluation chapter}{The evaluation chapter}{05-evaluation}
\deftplref{ch:related}{the related-work chapter}{The related-work chapter}{06-related-work}
\deftplref{fig:overview}{the leading figure}{The leading figure}{01-introduction}
\deftplref{fig:positioning}{the positioning figure}{The positioning figure}{06-related-work}
\deftplref{aidecl}{the AI declaration}{The AI declaration}{09-ai-declaration}

% A web link that looks like one: underlined, in the link color. Links set with \href or \url
% stay black, as the faculty template prints them (see \hypersetup above).
\definecolor{weblinkcolor}{RGB}{31,86,156}
\newcommand{\weblink}[2]{\href{#1}{\textcolor{weblinkcolor}{\uline{#2}}}}

% The Checks page of the draft: on every compile, the lint runs on the sources (shell escape, which
% Overleaf and the Makefile enable) and writes lint-checks.tex, the essentials and the other checks,
% input after the list of todos; its failing checks also come out as package warnings in the log.
% Without shell escape nothing is written and nothing is input. \lintfalse turns it off.
\newif\iflint\linttrue
\newcommand{\lintchecks}{\iflint\immediate\write18{texlua scripts/lint.lua overleaf --tex > lint-checks.tex}\InputIfFileExists{lint-checks.tex}{}{}\fi}

% Marks a step where the student stops for the advisors' approval (lifecycle chapter of the guide).
\newcommand{\anchor}{\textcolor{red!70!black}{\textsc{anchor}}}
\definecolor{official}{RGB}{255,243,205}
\definecolor{anchorrow}{RGB}{252,228,228} % lifecycle rows that are anchors
\definecolor{department}{RGB}{224,242,226} % lifecycle rows that are the department's own sessions
% What is yours, in every figure: the FTSRG TDK Howto asks to keep your own contribution visibly apart
% from prior work, and the leading figure is where the reviewer looks for it first. Use the `ours`
% TikZ style on the nodes that are your contribution and say so in the caption.
% ours: every component that is yours, one blue; problem red, result green; prior: reused or external, grey.
\tikzset{ours/.style={draw=blue!60!black, line width=1.1pt, fill=blue!8}, prior/.style={draw=gray!70, fill=gray!10},
  problem/.style={draw=red!60!black, fill=red!8}, result/.style={draw=green!50!black, fill=green!8}}
\newcommand{\department}{\textcolor{green!45!black}{\textsc{department}}}
\newcommand{\official}{\textcolor{orange!80!black}{\textsc{official}}}
\newcommand{\advisortodo}[1]{\todo[inline, color=cyan!20, bordercolor=cyan!60!black]{\textbf{Advisor:} #1}}

% Review comments from the supervisor in the student's paper:
%   \atodo[options]{comment}
% Pre-fills a green color and "A" (advisor) as author; like \todo, it prints inline unless
% given [noinline]. Every \todo option
% still works, and options given here override the defaults:
%   \atodo{review comment}  \atodo[noinline]{cite}  \atodo[author=IK]{...}
\newcommand{\atodo}[2][]{\todo[color=green!25, bordercolor=green!50!black, author=A, #1]{#2}}

% A gap the paper still has to fill, in a sentence, a table cell or a caption, where a margin note
% cannot go (a float has no margin):
%   \tbd{the proving time}          a value or a sentence still to write
%   \citeneeded  \citeneeded[the 2023 survey]   a claim that still needs its source
% Both print as coloured text in brackets in the draft, breaking across lines like the text around
% them, are listed with the todos, and disappear in the
% final build. Never type TBD, [CITE], [?] or XXX into the text: the markers lint fails on them,
% and on any \tbd or \citeneeded left at the submit phase.
\makeatletter
\newcommand{\tdk@gap}[2]{\if@todonotes@disabled\else{\color{#1}\sffamily\small[#2]}\addcontentsline{tdo}{todo}{#2}\fi}
\makeatother
\DeclareRobustCommand{\tbd}[1]{\csname tdk@gap\endcsname{orange!75!black}{TODO: #1}}
\DeclareRobustCommand{\citeneeded}[1][]{\if\relax\detokenize{#1}\relax\csname tdk@gap\endcsname{red!70!black}{citation needed}\else\csname tdk@gap\endcsname{red!70!black}{citation needed: #1}\fi}


%--------------------------------------------------------------------------------------
% Abbreviations, notation and terms (glossaries-extra)
%--------------------------------------------------------------------------------------
\setabbreviationstyle{long-short}          % first use "long form (SHORT)", then "SHORT"
\glssetcategoryattribute{abbreviation}{nohyperfirst}{true}
\renewcommand*{\glsnamefont}[1]{\textbf{#1}} % the first column in bold, as a table head would be
\setlength{\glsdescwidth}{0.50\linewidth}
\makenoidxglossaries
% \commonentry{key}{declaration} declares a common entry unless the paper has declared the key,
% and lists the key, so the guide can print the common glossary
\newcommand*{\tdkcommonkeys}{}
\newcommand{\commonentry}[2]{\ifglsentryexists{#1}{}{#2\listgadd{\tdkcommonkeys}{#1}}}
\loadglsentries{content/glossary}
\loadglsentries{include/glossary-common} % the entries our topics share; a paper's own declaration wins
\ifdefined\finalmode\else\loadglsentries{guide/glossary}\fi % the guide's own examples, only when the guide is built
% Set the section or definition that introduces an entry, including common entries.
% Call after the entries are loaded, e.g. in include/macros.tex:
% \glsintroduced{vc}{sec:credentials}
\newcommand{\glsintroduced}[2]{\glssetentryfield{#1}{user2}{\cref{#2}}}
\newglossarystyle{tdk-introduced}{%
  \setglossarystyle{long-booktabs}%
  \renewcommand*{\glossarysection}[2][]{\subsection*{##2}}%
  \renewenvironment{theglossary}{%
    \begin{longtable}{@{}p{0.18\linewidth}p{0.50\linewidth}p{0.22\linewidth}@{}}}{\end{longtable}}%
  \renewcommand*{\glossaryheader}{%
    \toprule\textbf{\entryname} & \textbf{\descriptionname} & \textbf{Introduced in}\\\midrule\endhead
    \bottomrule\endfoot}%
  \renewcommand*{\glossentry}[2]{%
    \glstarget{##1}{\glossentryname{##1}} & \glossentrydesc{##1} &
    \ifglshasfield{user2}{##1}{\glsentryuserii{##1}}{---}\tabularnewline}%
}
% The three lists, as booktabs tables with their own column heads, each with a sentence
% that says what the table is for. \printlists is called inside the appendix in content/10-appendices.tex.
% The appendix section, and each list in it, is left out while the text uses none of its entries (an entry
